\usepackage{babel}
\usepackage{amsmath,amssymb}
\usepackage{xspace}
\usepackage{amsthm}
\usepackage{physics}
\usepackage{mathtools}
\usepackage{latexsym}
\usepackage{stmaryrd}
\usepackage{microtype}
\usepackage{picins}

\usepackage{xparse}
\usepackage{environ} % for \Collect@Body
\usepackage[noadjust]{cite}

% `bigovee` symbol from mathabx fonts
\DeclareFontFamily{U}{mathux}{\hyphenchar\font45}
\DeclareFontShape{U}{mathux}{m}{n}{
      <5> <6> <7> <8> <9> <10>
      <10.95> <12> <14.4> <17.28> <20.74> <24.88>
      mathux10
      }{}
\DeclareSymbolFont{mathux}{U}{mathux}{m}{n}
\DeclareFontSubstitution{U}{mathux}{m}{n}
\DeclareMathSymbol{\bigovee}{1}{mathux}{"8F}

% Get `\coprod` symbol from Computer Modern
\DeclareSymbolFont{cmlargesymbols}{OMX}{cmex}{m}{n}
\DeclareMathSymbol{\cmcoprod}{\mathop}{cmlargesymbols}{"60}
\let\coprod\cmcoprod


\usepackage{xcolor}
\usepackage{tikz}
\usepackage{tikz-cd}
\usetikzlibrary{arrows, matrix}
\usepackage[textsize=footnotesize]{todonotes}


%\usepackage[inline]{showlabels}

\usepackage{enumitem}
\newlist{myenum}{enumerate}{2}
\setlist[myenum,1]{label=\textup{(\roman*)},ref=(\roman*)}
\setlist[myenum,2]{label=\textup{(\alph*)},ref=(\alph*)}

% fix eptcs.bst
\newcommand{\urlalt}[2]{\href{#1}{\urlstyle{rm}\nolinkurl{#2}}}

\newcounter{main}
\newtheorem{theorem}[main]{Theorem}
\newtheorem{proposition}[main]{Proposition}
\newtheorem{lemma}[main]{Lemma}
\newtheorem{corollary}[main]{Corollary}

\theoremstyle{definition}
\newtheorem{definition}[main]{Definition}
\newtheorem{notation}[main]{Notation}
\newtheorem{note}[main]{Note}
\newtheorem{example}[main]{Example}
\newtheorem{remark}[main]{Remark}

\DeclarePairedDelimiter{\ceil}{\lceil}{\rceil}
\DeclarePairedDelimiter{\floor}{\lfloor}{\rfloor}
\DeclarePairedDelimiter{\inn}{\langle}{\rangle}

\newcommand{\sse}{\subseteq}
\newcommand{\cl}[1]{\overline{#1}}
\DeclareMathOperator{\im}{im}
\newcommand{\asrt}{\mathrm{asrt}}
\newcommand{\id}{\mathrm{id}}
\newcommand{\opp}{\mathrm{op}}
\newcommand{\sa}{\mathrm{sa}}
\newcommand{\pred}{\mathrm{Pred}}
\newcommand\C{\mathbb{C}}
\newcommand\R{\mathbb{R}}
\newcommand\N{\mathbb{N}}
\newcommand\after{\mathbin{\circ}}
\DeclareMathOperator{\Idem}{Idem}
% \newcommand\TODO[1]{\textcolor{red}{\textbf{[}\emph{#1}\textbf{]}}}
\let\TODO\todo
\newcommand\Define[1]{\textbf{#1}}
\DeclarePairedDelimiter{\weight}{\lvert}{\rvert}
% commands added by Kenta
% % % % % % % % % % % % %

% hyphen in math mode
\DeclareMathSymbol{\mhyphen}{\mathalpha}{operators}{"2D}

\newcommand{\cat}[1]{\mathbf{#1}}
\newcommand{\catC}{\cat{C}}
\newcommand{\catD}{\cat{D}}

\newcommand{\calF}{\mathcal{F}}

\newcommand{\Rpos}{\mathbb{R}_{\ge 0}}
\newcommand{\Rpp}{\mathbb{R}_{> 0}}

\DeclareMathOperator{\absco}{absco}
\DeclareMathOperator{\co}{co}
\DeclareMathOperator{\Dom}{Dom}

% categories
\newcommand{\Set}{\mathbf{Set}}
\newcommand{\Pfn}{\mathbf{Pfn}}
\newcommand{\EA}{\mathbf{EA}}
\newcommand{\OA}{\mathbf{OA}}
\newcommand{\BA}{\mathbf{BA}}
\newcommand{\sOA}{\boldsymbol{\sigma}\mathbf{OA}}
\newcommand{\sBBNS}{\boldsymbol{\sigma}\mathbf{BBNS}}
\newcommand{\sBOUS}{\boldsymbol{\sigma}\mathbf{BOUS}}
%\newcommand{\Mod}{\mathbf{Mod}}
\newcommand{\CConv}{\mathbf{CConv}}
\newcommand{\Vect}{\mathbf{Vect}}
\newcommand{\OVSu}{\mathbf{OVSu}}
\newcommand{\OVSsu}{\mathbf{OVS}_\leq}
\newcommand{\OVSt}{\mathbf{OVSt}}
\newcommand{\SepConv}{\mathbf{SepConv}}
\newcommand{\pBNS}{\mathbf{pBNS}}
\newcommand{\OUS}{\mathbf{OUS}}
\newcommand{\BOUS}{\mathbf{BOUS}}
\newcommand{\SuperConv}{\mathbf{Conv}_{\infty}}
\newcommand{\CSuperConv}{\mathbf{CConv}_{\infty}}
%\newcommand{\sBpBNS}{\boldsymbol{\sigma}\mathbf{BpBNS}}

\newcommand{\Cstar}{\mathbf{Cstar}}
\newcommand{\Wstar}{\mathbf{Wstar}}

\newcommand{\pSet}{\mathbf{Set}_{*}}

\newcommand{\sEA}{\boldsymbol{\sigma}\mathbf{EA}}
\newcommand{\omegaBA}{\boldsymbol{\omega}\mathbf{BA}}

% M-EMod
% \newcommand{\MyNewModLikeCat}[2]{%
%   \NewDocumentCommand{#1}{o}{\IfValueT{##1}{##1\mhyphen}#2}}
% EMod_M
\newcommand{\MyNewModLikeCat}[2]{%
  \NewDocumentCommand{#1}{o}{#2\IfValueT{##1}{_{##1}}}}
\MyNewModLikeCat{\WMod}{\mathbf{WMod}}
\MyNewModLikeCat{\CWMod}{\mathbf{CWMod}}
\MyNewModLikeCat{\sWMod}{\boldsymbol{\sigma}\mathbf{WMod}}
\MyNewModLikeCat{\sCWMod}{\boldsymbol{\sigma}\mathbf{CWMod}}
\MyNewModLikeCat{\sEMod}{\boldsymbol{\sigma}\mathbf{EMod}}
\MyNewModLikeCat{\EMod}{\mathbf{EMod}}
\MyNewModLikeCat{\AEMod}{\mathbf{AEMod}}
\MyNewModLikeCat{\Conv}{\mathbf{Conv}}


\newcommand{\Kl}{\mathcal{K}\mspace{-2mu}\ell}
\newcommand{\EM}{\mathcal{E}\mspace{-3mu}\mathcal{M}}


% functors
\newcommand{\Dst}{\mathcal{D}}
\newcommand{\Pow}{\mathcal{P}}
\newcommand{\iDst}{\Dst^{\infty}}

\newcommand{\truth}{\mathbf{1}}
\newcommand{\falsity}{\mathbf{0}}
\newcommand{\mpunct}[1]{\,\text{#1}}

\newcommand{\Tot}{\mathrm{Tot}}
\newcommand{\St}{\mathrm{St}}
\newcommand{\sSt}{\St_{\le}}
\newcommand{\Pred}{\mathrm{Pred}}
\newcommand{\Base}{B}
\newcommand{\sBase}{\Base_{\le}}
\newcommand{\pproj}{\mathord{\vartriangleright}}
\newcommand{\copow}{\cdot}

\newcommand{\Tz}{\mathcal{T}} % Totalization
\newcommand{\Mlt}{\mathcal{M}}

\newcommand{\smul}{\cdot}

\DeclarePairedDelimiter{\paren}{(}{)}
\DeclarePairedDelimiter{\tup}{\langle}{\rangle}
\DeclarePairedDelimiter{\cotup}{[}{]}
\DeclarePairedDelimiter{\ptup}{\langle\!\langle}{\rangle\!\rangle}

\newcommand{\Comment}[1]{%
  \marginpar{\raggedright\hbadness=10000
    \def\baselinestretch{0.8}\footnotesize
    #1\par}}

% \set command
% eg. \set[\Big]{x\in X | f(x) = y}
\makeatletter
\newcommand{\my@set@mid}{\mathrel{}\mathclose{}\delimsize\vert\mathopen{}\mathrel{}}
% see https://tex.stackexchange.com/questions/25398/mathrel-conflict-with-left-right-middle
\DeclarePairedDelimiterX{\set}[1]{\lbrace}{\rbrace}{%
  \my@repl@vert@i{\my@set@mid}{#1}}
% Note: \DeclarePairedDelimiterX creates group
\NewDocumentCommand{\my@repl@vert@i}{m >{\SplitArgument{1}{|}}m}{\my@repl@vert@i@{#1}#2}
\NewDocumentCommand{\my@repl@vert@i@}{m m m}{\IfNoValueTF{#3}{#2}{#2#1#3}}
\makeatother

% align with \textstyle
% https://tex.stackexchange.com/questions/274449/is-there-an-environment-that-imitates-align-but-displays-textstyle-math
\newenvironment{talign}
 {\let\displaystyle\textstyle\align}
 {\endalign}
\newenvironment{talign*}
 {\let\displaystyle\textstyle\csname align*\endcsname}
 {\endalign}
% for amsthm
\makeatletter
\let\talign\align@qed
\expandafter\let\csname talign*@qed\endcsname\align@qed
\makeatother

% Auxproof
\newcommand{\auxprooffont}{\color{black!60}}
\newenvironment{Auxproof}{}{}

\newcommand{\showauxproof}{%
  \renewenvironment{Auxproof}
  {\par
    \auxprooffont\noindent\textbf{AUX-PROOF}\dotfill\par
    \noindent\ignorespaces}
  {\par\noindent\textbf{AUX-QED}\dotfill\par
    \noindent\ignorespacesafterend}%
}
\makeatletter
\newcommand{\hideauxproof}{%
  \renewenvironment{Auxproof}{\@bsphack\Collect@Body\@gobble}{\@Esphack}%
}
\hideauxproof
\makeatother
